\appendix

\chapter{Complete Run Registry}
\label{app:registry}
Table~\ref{tab:registry} lists the runs behind the Phase-1 findings. All runs log to
Weights \& Biases (project \texttt{rlvr-openings}) with per-step telemetry and saved
group reward tensors; runs are pinned by config, seed, and framework.

\begin{table}[h]
\centering\small
\begin{tabular}{@{}L{0.13\linewidth}L{0.32\linewidth}L{0.14\linewidth}L{0.28\linewidth}@{}}
\toprule
\textbf{Study} & \textbf{Model / setup} & \textbf{Seeds} & \textbf{Primary output} \\
\midrule
P1 (toy) & Qwen2.5-1.5B, 28 layers, 5 steps, arithmetic & 1 & per-layer grad-norm; overlap$=1.0$ \\
P1 (scaled) & Qwen2.5-3B, 36 layers, 10 steps, GSM8K & $\{0,1,2,3\}$ & overlap $0.083\pm0.048$; top-25\% share $0.39$ \\
P2 & 4 methods, 40 steps, $16\times8$, GSM8K & shared seed-0 sched. & ZVF $0.72$--$0.77$ \\
P3 & Qwen3.5-4B, $G\in\{2,4,8,16\}$, 6 steps & 0 & token cost $9.7$k--$68$k; exploratory sweep \\
Curriculum & Qwen3.5-4B, 8 steps, mixed-variance filter & 0 & ZVF $0.50\to0.00$, $4.81\times$ cost \\
Matched-token & Qwen3.5-4B, $G{=}4$, baseline vs curriculum, ${\sim}30$k tokens & $\{0,1,2\}$ & both arms $+0.028$ mean \\
P8 & method-trace classifier, $N=160$ correlated rows, 5-fold row CV & 1 trace/method & AUROC $0.838\pm0.010$ \\
\bottomrule
\end{tabular}
\caption{Complete Phase-1 run registry.}
\label{tab:registry}
\end{table}

\chapter{Reproducibility and Provenance}
\label{app:repro}
Every published diagnostic in this report is regenerable from stored artefacts:
\begin{itemize}[leftmargin=1.4em,itemsep=1pt]
  \item \texttt{p2\_collapse\_analysis.py} regenerates the ZVF / collapse JSON from
  the saved P2 reward tensors (reproduces $\zvf=0.72$--$0.77$).
  \item \texttt{p8\_detector.py} regenerates \texttt{detector\_reproduced.json} from
  the \path{n2\_reward\_tensor\_resume} tensors (reproduces row-CV AUROC $0.838$).
  \item \path{experiments/results/token\_budget/results.json} records the matched-token
  three-seed baseline-vs-curriculum comparison.
\end{itemize}
Platforms and budget: managed GRPO training via Tinker; white-box P1 on a Google
Colab L4 GPU; batch scoring on Modal; logging on Weights \& Biases. The full
Phase-1 empirical portfolio was executed within a modest budget (order of tens of
US dollars of managed compute plus free-tier GPU time).

\paragraph{Adversarial review and verification.} AI-assisted adversarial review was
used during development to identify overclaims and prompt additional checks. These
model assessments are not treated as scientific evidence and are not required to
reproduce any number in the report. Verification in the submission is script-based,
using stored tensors and explicit result JSONs as described above.

\paragraph{AI-assistance disclosure.} Language models assisted with adversarial
review, code review, and editorial revision. The candidate remains responsible for
the research design, execution choices, source verification, calculations, and final
claims. No model-generated verdict is reported as an empirical result.

\paragraph{Code availability.} All harness code, diagnostic scripts, raw reward tensors,
and result JSONs are maintained in the project repository under version control, pinned
by configuration, seed, and framework; the run registry (Table~\ref{tab:registry}) keys
each finding to its producing run. The complete report source, figures, and this
document are likewise version-controlled.

\paragraph{Originality / plagiarism.} This report is an original individual work. All
external sources and prior art are cited (Chapter~\ref{ch:lit}, References). The report
has been prepared for the institutional similarity (plagiarism) check with a target
similarity of $\le 15\%$ in line with the programme guidelines, and reuses the
candidate's own prior project write-ups where applicable.

\chapter{Notation, Symbols, and Abbreviations}
\label{app:glossary}
\small
\begin{longtable}{@{}L{0.22\linewidth}L{0.70\linewidth}@{}}
\toprule
\textbf{Term / symbol} & \textbf{Meaning} \\
\midrule
\endhead
GRPO & Group Relative Policy Optimization; the group-baseline RL post-training
algorithm studied here~\cite{shao2024deepseekmath}. \\
RLVR & Reinforcement Learning with Verifiable Rewards; a programmatic checker
supplies the reward with no learned reward model. \\
ZVF ($\zvf$) & Zero-Variance Fraction; the fraction of prompt-groups in a step whose
within-group reward variance is zero (Eq.~\eqref{eq:zvf}). \\
$G$ & Group size; number of sampled completions drawn per prompt. \\
$A_i$ & Group-relative advantage of completion $i$ (Eq.~\eqref{eq:adv}). \\
$\mu_g,\sigma_g$ & Within-group reward mean and standard deviation. \\
$\beta$ & KL-penalty coefficient toward the frozen reference policy. \\
Gradient utilisation & Complementary fraction of groups with non-zero reward-relative
advantage signal ($1-\zvf$ at the group level); excludes KL-only updates. \\
Held-out gain & Change in accuracy on a held-out prompt set; always reported with
its denominator $n$. \\
All-correct / all-wrong & Decomposition of collapsed (zero-variance) groups by
whether every completion was right or wrong. \\
Matched baseline & A control run identical to the treatment except for the single
lever under study. \\
AUROC & Area under the ROC curve; the exploratory P8 method-trace classifier metric. \\
\bottomrule
\end{longtable}
\normalsize
